\label{sec-related}
\section{Related Work}
\subsection{Forward Referencing as a Lure} 
% \paragraph{Forward Referencing as a Lure}
%Our paper contributes to a growing body of work on developing natural language generation models from a psychological aspect. 
%Several works in the journalism and psychology field have investigated how clickbait headlines influence readers reading preference. 
\citet{psychologyofcuriosity} shows how the desire for information motivates
human curiosity. 
%To attract clicks and views, curiosity gaps are also used in online media when composing clickbait-style titles.
Forward-referencing has been defined as a technique for creating curiosity gaps
at a discourse level for use in headlines~\cite{BLOM201587,Yang2011ACI}. 
A similar concept is cataphora, in which information is forwarded as a teaser
at the sentence level~\cite{cataphoric, Halliday76a}. 
\citet{digitalclickbait} investigate how editors rewrite headlines for digital
platforms, and analyze the linguistic features of what makes for an attractive
headline.
\citet{zhang2018question} address attractive headline generation as question
headline generation (QHG), which assumes that interrogative sentences are more
popular. Although this modality is indeed a type of FR, we argue that the
interrogative style may not be suitable for all kinds of headlines, especially
verified news. Hence in our work, we fully consider all kinds of FR which
are commonly used and seen in social media and on digital platforms.
% We summarize those FRs as below:
% \begin{itemize}
%     \item \textbf{Type~0}: None of the below
%     \item \textbf{Type~1}: Demonstrative pronouns
%     \item \textbf{Type~2}: Personal pronouns
%     \item \textbf{Type~3}: Definite articles
%     \item \textbf{Type~4}: Ellipses
%     \item \textbf{Type~5}: Imperatives
%     \item \textbf{Type~6}: Interrogatives
%     \item \textbf{Type~7}: General nouns
%     \item \textbf{Type~8}: Location Adverbs
% \end{itemize}
%Table~\ref{forward-ref-tab} summarizes the types of forward references, and 
Sample headlines exhibiting FR techniques can be found in
Fig.~\ref{fr-demo} in the appendix.

% \begin{table}[h]
% \centering
% \small
% \caption{\label{font-table} Types of forward references} 
% \begin{tabular}{ll}
% \hline {}&\textbf{Forward-reference Types} \\ 
% \hline
% Type~0 & None of the below\\
% Type~1 & Demonstrative pronouns\\
% Type~2 & Personal pronouns\\ 
% Type~3 & Definite articles\\ 
% Type~4 & Ellipses\\
% Type~5 & Imperatives\\
% Type~6 & Interrogatives\\
% Type~7 & General nouns\\
% Type~8 & Location Adverbs\\
% \hline
% \end{tabular}
% \label{forward-ref-tab}
% \end{table}
%We utilize different FR as a hint to generate more general and appealing headlines.
%Recent research \cite{digitalclickbait} also investigated the process of how editors re-write the news headline for digital platforms and analyzed the linguistic feature of what a effective headline contains. This work provide us a deeper understanding of the process of transforming news headlines to attractive or clickbait-style.
\subsection{Headline Generation} 
%\paragraph{Headline Generation}
Headline generation can be viewed as a more specific summarization task.
\citet{qi2020prophetnet} propose a Transformer-based, self-supervised n-gram
prediction objective. 
\citet{liu2019finetune} propose BERTSum, a variation of
BERT~\cite{devlin-etal-2019-bert} for extractive summarization.
%\citet{2020t5} propose a text-to-text Transformer-based generation model.
\citet{see-etal-2017-get} propose an attention-based pointer generator with a
copy mechanism, which has made great progress in summarization. 
Although its ability to copy text from the source context is powerful,
using it directly for verified news often leads to bland titles.
Hence we apply FRs and a sensationalism scorer to produce more satisfying results.
\citet{xu-etal-2019-clickbait} propose auto-tuned reinforcement learning
to generate sensational headlines using a pretrained sensationalism scorer;
the resulting score is used as the reward to enhance the attractiveness.
Although generating attractive headlines has been widely
explored~\cite{popularitybased, jin-etal-2020-hooks}, we focus more on
fidelity to ensure that the semantics of the generated headline are faithful to
the source content to avoid harmful hallucination.

\subsection{Faithful Summarization}
% \paragraph{Faithful Summarization}
Recent work investigates how to improve the faithfulness of the generated
summary or headline. \citet{Matsumaru2020ImprovingTO} propose 
% an intrinsic method by 
pretraining a textual entailment scorer to filter out noisy samples
in the dataset, preventing hallucination or unfaithful generation.
\citet{onfaith} analyze the faithfulness of current abstractive summarization
systems, and discover that textual entailment is better correlated
to faithfulness than standard metrics. Based on such work, one
major direction is to evaluate generated summaries in terms of textual entailment
rather than raw metrics such as ROUGE~\cite{lin-2004-rouge} or BLEU~\cite{bleu}.
Accordingly, we propose a faithfulness scorer based on textual entailment to
evaluate how well the generated headlines fit the semantics of the content. 
